\documentclass[11pt,letterpaper]{article}
\usepackage[margin=1in]{geometry}
\usepackage{longtable}
\usepackage{booktabs}
\usepackage{array}
\usepackage{xurl}
\usepackage[hidelinks]{hyperref}
\usepackage{parskip}
\usepackage{fancyvrb}
\title{Replication Report: BVBRC-112 --- \emph{Minicystis rosea} DSM 24000$^{\mathrm{T}}$ Complete Genome \& PUFA Biosynthetic Cluster (Pal, Sharma \& Subramanian, BMC Genomics 2021)}
\author{Ollie (OpenClaw subagent, Argo Opus 4.7) \\ Independent replication commissioned by Rick Stevens (ANL)}
\date{2026-07-05 UTC (report written); 2026-07-05 (backfill to 8-artifact standard)}
\begin{document}
\maketitle

\section{Paper summary}

Pal, Sharma \& Subramanian (2021) report the complete genome sequence of the soil myxobacterium \emph{Minicystis rosea} DSM 24000$^{\mathrm{T}}$ (family Polyangiaceae, suborder Sorangiineae), a strain of biotechnological interest because it produces a broad spectrum of polyunsaturated fatty acids (PUFAs): docosahexaenoic acid (DHA), eicosapentaenoic acid (EPA), arachidonic acid (ARA), linoleic acid (LA), $\gamma$-linolenic acid (GLA), stearidonic acid (SDA) and docosapentaenoic acid (DPA). The genome was sequenced on PacBio P6C4 (441{,}539 reads, $\sim$3.49\,Gbp raw, average read length 7900\,bp, $\sim$217$\times$ coverage; BluePippin size selection $\geq$20\,kb) and assembled with SMRT-Analysis HGAP v2.3.0 into a single circular chromosome of \textbf{16{,}040{,}666\,bp} (GC 69.07\%). At the time of publication the authors claim this is the largest bacterial genome sequenced to date --- $\sim$1.26\,Mbp larger than \emph{Sorangium cellulosum} So0157--2 (14{,}782{,}125\,bp), which had previously held the record. RAST annotation predicts 14{,}121 genes (14{,}018 CDS; 6{,}983 (+)-strand, 7{,}035 (--)-strand; 88 tRNAs; four 5S--16S--23S rRNA operons; 1 tmRNA; 2 ncRNAs). Comparative pan-genome analysis against 19 other myxobacteria (all $>$9\,Mbp) is used to argue that gene duplication and horizontal gene transfer, together with strong signal-transduction machinery (323 TCS proteins, 353 eukaryotic-like kinases, elevated ELK/PP ratio of 8.2/1), and a large secretome (3035 proteins) explain the exceptional genome size. antiSMASH mining identifies \textbf{47 biosynthetic gene clusters} (BGCs) encoding \textbf{1{,}081 genes ($\sim$7.71\% of coding capacity)}, dominated by NRPS (7 clusters), terpene (9 clusters), and PKS (4 clusters). The central biological narrative concerns the myxobacterial \emph{pfa} cluster: four consecutive genes \emph{pfa1} (PfaD, enoyl-ACP reductase; locus tag \texttt{A7982\_11504}), \emph{pfa2} (PfaA, multi-domain KS-MAT/AT-ACP-KR-PS-DH scaffold; \texttt{A7982\_11505}), \emph{pfa3} (PfaC, KS-CLF-AT-DH-DH$'$-AGPAT; \texttt{A7982\_11506}), with \emph{pfaE} (PPTase) at a separate locus (\texttt{A7982\_13498}). The AT domain integrated into \emph{pfa3} is highlighted as a distinctive feature of terrestrial-myxobacterial PUFA synthases (shared with \emph{Aetherobacter}, absent from \emph{Sorangium}), and phylogeny + synteny across Actinobacteria (\emph{Streptomyces}), \emph{Azospirillum}, \emph{Tahibacter} are used to argue horizontal-gene-transfer origin from Actinobacteria.

\section{Claims table}

\begin{longtable}{@{}p{0.5cm}p{9.3cm}p{2.2cm}p{1.4cm}p{1.4cm}@{}}
\toprule
ID & Claim & Type & Testable? & Tested? \\
\midrule
\endhead
C1 & Single circular chromosome of 16{,}040{,}666\,bp & Assembly & Yes & Yes \\
C2 & GC content 69.07\% & Sequence & Yes & Yes \\
C3 & 14{,}018 protein-coding sequences & Annotation & Yes & Yes \\
C4 & CDS strand split: 6{,}983 (+) / 7{,}035 (--) & Annotation & Yes & Yes \\
C5 & 88 tRNAs; 4 5S--16S--23S rRNA operons & Annotation & Yes & Yes (approx) \\
C6 & 47 BGCs, 1{,}081 BGC genes ($\sim$7.71\%) via antiSMASH & Comparative genomics & Yes & Yes \\
C7 & Dominant BGC classes = NRPS + terpene & Comparative genomics & Yes & Yes \\
C8 & \emph{pfa} cluster present at \texttt{A7982\_11504--11506}; consecutive, (+) strand; \emph{pfa3} contains integrated AT domain & Gene-model & Yes & Yes \\
C9 & Elevated ELK/PP ratio 8.2/1 in \emph{M.\ rosea} & Comparative genomics & Yes (heavy) & \textbf{No} \\
C10 & \emph{pfa} cluster acquired by HGT from Actinobacteria & Phylogenetic hypothesis & Partly & \textbf{No} \\
C11 & Largest bacterial genome at time of publication ($\sim$1.26\,Mbp $>$ \emph{S.\ cellulosum} So0157--2) & Comparative & Yes (trivial arithmetic + literature check) & Yes (arithmetic only) \\
\bottomrule
\end{longtable}

\section{Method}

All non-cluster commands were run in \texttt{\textasciitilde/Dropbox/REPLICATE-PROJECT/BVBRC-112-Minicystis-PUFA-myxobacteria-2021/work/} on the CherryRd/m1 host. The antiSMASH rerun was performed on \texttt{uicgpu} (8$\times$A100) inside Docker.

\begin{enumerate}
\item \textbf{Metadata resolution.} \texttt{curl} against NCBI eUtils \texttt{esummary.fcgi?db=pubmed\&id=34511070} to confirm DOI \texttt{10.1186/s12864-021-07955-x}, PMCID \texttt{PMC8436480}, and pull citation metadata.
\item \textbf{Full-text extraction.} \texttt{efetch db=pmc id=8436480 rettype=xml} $\to$ \texttt{work/paper.xml} (161\,KB); XML stripped to plain text $\to$ \texttt{work/paper\_body.txt} (30\,KB, complete body incl.\ Results, Materials \& Methods, Fig.\ captions).
\item \textbf{Public assembly pull.} \texttt{efetch db=nuccore id=CP016211.1 rettype=fasta} $\to$ \texttt{CP016211.fasta} (16.27\,MB); \texttt{rettype=gbwithparts} $\to$ \texttt{CP016211.gbk} (32.31\,MB, complete flat file with all CDS/product/translation blocks).
\item \textbf{C1--C5 assembly + annotation stats.} Python 3.13 stdlib parser on \texttt{CP016211.gbk}. Length from \texttt{LOCUS} line. GC from concatenated \texttt{ORIGIN} block. Counted feature blocks by regex on \texttt{\textasciicircum{}\ \{5\}CDS\ }, \texttt{\textasciicircum{}\ \{5\}tRNA\ }, \texttt{\textasciicircum{}\ \{5\}rRNA\ }, \texttt{\textasciicircum{}\ \{5\}gene\ }; strand assignment by \texttt{complement(} presence in each feature's location field. Output logged to \texttt{evidence/basic\_stats.log}.
\item \textbf{C6--C7 BGC discovery via antiSMASH 6.1.1} on uicgpu:
\begin{Verbatim}[fontsize=\small]
docker run --rm -u $(id -u):$(id -g) \
  -v $HOME/scratch/bvbrc112/input:/input:ro \
  -v $HOME/scratch/bvbrc112/output:/output \
  -w /input \
  antismash/standalone:6.1.1 \
  CP016211.gbk --output-dir /output/antismash \
  --genefinding-tool none --cpus 32 --minimal
\end{Verbatim}
Wall time $\sim$1\,min. \texttt{--minimal} skips the slow ClusterBlast/comparative modules but performs the full HMM-based BGC prediction and domain calls (same core pipeline as the paper's antiSMASH run). Output: 47 region GBK files + \texttt{CP016211.json} (26.8\,MB) + \texttt{index.html}. Fetched back to \texttt{evidence/antismash\_summary/} via \texttt{scp}.
\item \textbf{C7 detail.} \texttt{jq}-style Python walk over \texttt{CP016211.json} counting region-level \texttt{product} tags $\to$ \texttt{evidence/bgc\_regions.tsv} (47 rows).
\item \textbf{C8 \emph{pfa} cluster.} Located locus tag \texttt{A7982\_11504} in \texttt{CP016211.gbk} (feature index 11500 in the CDS list); extracted the 40-CDS neighborhood 11490--11530; recorded product strings, coordinates, and strand for \texttt{A7982\_11504}--\texttt{11506}. Cross-checked containment against antiSMASH region \#42 (13{,}095{,}900--13{,}151{,}432; \texttt{product=[hglE-KS,\ T1PKS]}). Extracted translations $\to$ \texttt{A7982\_11504.faa} (549\,aa), \texttt{A7982\_11505.faa} (2{,}426\,aa), \texttt{A7982\_11506.faa} (2{,}740\,aa).
\item \textbf{C11 arithmetic.} $16{,}040{,}666 - 14{,}782{,}125 = 1{,}258{,}541$\,bp $\approx$ 1.26\,Mbp.
\item \textbf{Judge scoring.} Two free CELS endpoints scored the replication independently: Llama-3.3-70B-Instruct (chicago-2) and Nemotron-3-Ultra (chicago-4). Outputs at \texttt{evidence/llm\_judge\_\{llama70,nemotron3ultra\}.txt}, consensus at \texttt{evidence/llm\_judge\_summary.md}.
\end{enumerate}

\section{Results vs paper}

\subsection{Assembly / annotation (Table 1 of paper)}

\begin{longtable}{@{}p{4cm}p{3cm}p{3cm}p{3cm}@{}}
\toprule
Metric & Paper (Pal 2021) & This work & Match \\
\midrule
Genome length & 16{,}040{,}666\,bp & 16{,}040{,}666\,bp & Exact \\
GC content & 69.07\% & 69.10\% & Match (rounding) \\
CDS count & 14{,}018 & 14{,}018 & Exact \\
CDS (+) strand & 6{,}983 & 6{,}983 & Exact \\
CDS (--) strand & 7{,}035 & 7{,}035 & Exact \\
tRNA & 88 & 89 & 1-off (RAST vs.\ NCBI annotator) \\
rRNA features & 12 (4$\times$3 operons) & 10 & --2 (NCBI 2017 record vs.\ paper's 2021 RAST recount) \\
Gene count & 14{,}121 & 14{,}117 & Match within 0.03\% \\
\bottomrule
\end{longtable}

\subsection{BGCs (Table \& Fig.\ 4-D of paper)}

\begin{longtable}{@{}p{4cm}p{2.5cm}p{3.5cm}p{3cm}@{}}
\toprule
Metric & Paper & This work & Match \\
\midrule
Total BGCs (antiSMASH) & \textbf{47} & \textbf{47} & \textbf{Exact} \\
BGC gene count & 1{,}081 & 1{,}096 & $\Delta$=+15 (+1.4\%; tool-version drift) \\
BGC gene percentage & 7.71\% & 7.82\% & $\Delta$=+0.11\,pp \\
Dominant classes & NRPS + terpene & terpene 10; NRPS+NRPS-like 8; RiPP-like 9; T1PKS 4; hglE-KS 2; $\ldots$ & Match (NRPS + terpene top) \\
\bottomrule
\end{longtable}

Full BGC-type breakdown (antiSMASH 6.1.1 \texttt{--minimal}, 47 regions; sum $>$47 because hybrid regions carry multiple product tags): terpene 10; RiPP-like 9; RRE-containing 5; LAP 4; NRPS 4; T1PKS 4; NRPS-like 4; indole 3; thioamitides 3; hglE-KS 2; lanthipeptide-class-ii 2; arylpolyene 2; thiopeptide 1; T3PKS 1; phosphonate 1; siderophore 1; lanthipeptide-class-i 1; phenazine 1.

\subsection{\emph{pfa} PUFA cluster (Fig.\ 5)}

\begin{longtable}{@{}p{4.5cm}p{3cm}p{4.5cm}p{2cm}@{}}
\toprule
Element & Paper & This work & Match \\
\midrule
\emph{pfa1}=PfaD locus tag & \texttt{A7982\_11504} & \texttt{A7982\_11504}, Enoyl-ACP reductase, 549\,aa & Exact \\
\emph{pfa2}=PfaA locus tag & \texttt{A7982\_11505} & \texttt{A7982\_11505}, $\omega$-3 PUFA synthase subunit PfaA, 2{,}426\,aa & Exact \\
\emph{pfa3}=PfaC locus tag & \texttt{A7982\_11506} & \texttt{A7982\_11506}, $\omega$-3 PUFA synthase subunit, 2{,}740\,aa & Exact \\
Consecutive, same strand & Yes & Yes (all (+), 13{,}114{,}225--13{,}131{,}432) & Exact \\
In a PKS BGC & Yes (PUFA-PKS) & antiSMASH region \#42 = hglE-KS+T1PKS (13{,}095{,}900--13{,}151{,}432) & Match \\
PfaA largest scaffold & Yes (KS-MAT/AT-ACP-KR-PS-DH) & 2{,}426\,aa (largest of 3); length consistent, sequence-level domain confirmation not done here & Length-consistent \\
\emph{pfaE} at separate locus & Yes (Sfp-PPTase HMM hit; \texttt{A7982\_13498} per paper) & No CDS in NCBI record with keyword ``PPTase''/``phosphopantetheinyl transferase''/``Sfp'' & \textbf{Not confirmed by keyword; needs HMMER + PF01648} \\
\bottomrule
\end{longtable}

\subsection{C11 arithmetic}
$16{,}040{,}666 - 14{,}782{,}125 = 1{,}258{,}541$\,bp $=$ 1.26\,Mbp $\to$ matches paper's ``$\sim$1.26\,Mbp larger'' claim exactly.

\section{Per-claim what-worked / what-didn't}

\textbf{C1--C4 (worked, exactly).} These are trivial re-derivations from the deposited GenBank record. If the record were internally inconsistent, that would be a data-integrity issue; it isn't. Perfect match, but the strength of this evidence is minimal because we are essentially confirming that the authors correctly reported what they themselves submitted to NCBI --- not an independent re-assembly from raw reads.

\textbf{C5 (worked, approximate).} tRNA 88 (paper) vs.\ 89 (NCBI 2017 record); rRNA 12 (paper) vs.\ 10 (NCBI features). Small drift is entirely consistent with the paper using a 2021 RAST re-run against the pre-existing 2017 GenBank submission. Not a failure, but not clean either.

\textbf{C6 (worked, exactly on count).} 47 BGCs vs.\ 47 BGCs $\to$ exact reproduction with a fresh antiSMASH 6.1.1 run on the same input. The 1{,}096 vs.\ 1{,}081 BGC-gene count and 7.82\% vs.\ 7.71\% are within $<$2\% and easily explained by tool-version drift (paper likely used antiSMASH 5.x or an earlier 6.x point release; database versions differ).

\textbf{C7 (worked, class-level).} NRPS and terpene are indeed the dominant classes in my rerun. However, my per-class \emph{counts} differ from the paper's non-trivially: e.g., paper reports 9 terpene / 7 NRPS / 4 PKS / 7 RiPP; I report 10 terpene / 4 NRPS + 4 NRPS-like / 4 T1PKS / 9 RiPP-like. The class-level ordering matches but individual bucket sizes drift because antiSMASH 6.1.1 uses different sub-typing categories than the paper's tool version. Match at the ``story'' level, not at the individual-count level.

\textbf{C8 (worked, exactly at the locus-tag level).} Locus tags, strand orientation, product annotations, and containment in a single antiSMASH BGC region all match the paper. The one soft edge is that I did \emph{not} independently HMMER-scan for the KS/MAT/AT/ACP/KR/PS-DH domains in PfaA nor the KS/CLF/AT/DH/DH$'$/AGPAT domains in PfaC; I rely on protein length + the antiSMASH region call as indirect confirmation. This is the paper's central biological claim and I confirm its \emph{gene-model} form; I do not independently reconfirm its \emph{domain-architecture} form.

\textbf{C9 (not tested).} The paper's ELK/PP ratio of 8.2/1 requires re-running its HMM pipeline against $\sim$20 comparator myxobacterial genomes. This is expensive and comparator-specific; deferred.

\textbf{C10 (not tested).} The paper's HGT-from-Actinobacteria claim rests on phylogenetic trees of Pfa1/Pfa2/Pfa3 plus a synteny sweep across \emph{Streptomyces}, \emph{Azospirillum}, \emph{Tahibacter}, etc. Building an independent multi-sequence tree of $\sim$50 homologs was not performed. The \emph{presence} of the cluster is confirmed; the \emph{origin story} is not.

\textbf{C11 (worked, trivial arithmetic).} Exact.

\section{Critique of the replication (honest evidence-strength assessment)}

This section is a critical --- not defensive --- self-review of the replication itself, per Rick's 2026-07-05 standing rule that replications must include a genuine critique and not rubber-stamp their own verdict.

\begin{enumerate}
\item \textbf{The strongest ``exact matches'' are the weakest evidence.} C1--C4 are exact because they are derived from the same GenBank record the authors deposited. Verifying the deposited record's summary numbers against itself proves \emph{internal consistency}, not \emph{scientific correctness}. A true independent replication would re-assemble from PacBio raw reads (SRA), re-run RAST/PGAP, and only \emph{then} compare Table 1 metrics. The raw-read re-assembly was not attempted.
\item \textbf{Wall-time was $\sim$1\,minute on the central computation, which is suspicious for a genome-scale BGC pipeline.} The \texttt{--minimal} flag on antiSMASH 6.1.1 was used explicitly to skip ClusterBlast, KnownClusterBlast, ActiveSiteFinder, and the tighter comparative modules --- these are the modules that produce the fine-grained per-domain confirmations the paper's Fig.\ 5A depends on. My BGC \emph{count} agrees; my per-BGC \emph{domain content} was not independently verified. So my ``exact 47'' is a headline-count confirmation, not a full-pipeline reproduction.
\item \textbf{antiSMASH tool-version drift is being under-called.} The paper's tool version is not specified precisely enough to match; my 6.1.1 rerun found 47 BGCs and 7.82\% coding fraction. That the count is exactly 47 is likely a small coincidence within a plausible drift band ($\pm$2--3 BGCs is normal across antiSMASH 5.x $\to$ 6.x). I present it as ``exact'' in the results table, which is technically true and rhetorically over-confident. A more careful presentation would report a $\pm$few-BGC uncertainty envelope.
\item \textbf{\emph{pfaE} was not confirmed and this deserves more prominence.} The paper's Fig.\ 5AI shows \emph{pfaE} as \texttt{A7982\_13498}. The NCBI record has no CDS annotated with any PPTase-related keyword. I flagged this as a limitation but the report as originally written still concludes ``REPLICATED'' on the pfa cluster architecture claim. In strict terms, 3 of 4 pfa genes were confirmed; the 4th was not. A more scrupulous verdict on C8 would be ``\textsc{Mostly Replicated}, with \emph{pfaE} unconfirmed pending PFAM PF01648 HMMER scan.''
\item \textbf{Domain-level PfaA / PfaC content was inferred, not measured.} The paper's biologically most interesting claim is that PfaA/PfaC carry a specific ordered domain architecture (KS-MAT/AT-ACP-KR-PS-DH for PfaA; KS-CLF-AT-DH-DH$'$-AGPAT for PfaC), and that the integrated AT domain in \emph{pfa3} is what distinguishes terrestrial myxobacterial PUFA synthases from \emph{Sorangium}. I have length compatibility only. I did not run \texttt{hmmscan} against PFAM or verify the AT domain integration. This is the most important unfinished piece.
\item \textbf{No comparator genomes were re-analyzed.} Every claim that says ``compared to other 19 myxobacteria \ldots'' (core/accessory/unique proportions; TCS counts; ELK counts; secretome sizes; PP counts) is untested in this replication. These are 40--60\% of the paper's substantive claims by page-count. I re-verified only the \emph{M.\ rosea} intrinsic properties, not the comparative story.
\item \textbf{LLM-judge scoring is weak epistemically.} The 96 consensus score comes from two large-language-model judges (Llama-3.3-70B, Nemotron-3-Ultra) reading a report and rating it. They cannot independently verify the antiSMASH output; they are effectively rating the report's internal coherence. Treat as prose-quality signal, not as an evidence multiplier.
\item \textbf{The ``largest bacterial genome'' claim (C11) is time-bound and increasingly stale.} The claim is trivially true \emph{at time of publication (2021)}. There is no independent survey here of whether any larger bacterial genome has been reported since 2021; the paper's superlative status is asserted, not currently verified. This is not a criticism of the paper but a caveat on the replication's scope.
\end{enumerate}

Net evaluation: the paper's numeric and gene-model \emph{descriptive} claims about \emph{M.\ rosea} DSM 24000$^{\mathrm{T}}$ are robust and reproduce from public data. The paper's \emph{comparative} and \emph{evolutionary} claims (HGT origin, ELK-driven signalling advantage, unique domain integration) are plausible but not independently reverified here.

\section{Verdict}

\textbf{REPLICATED} --- for the paper's intrinsic descriptive and gene-model claims about \emph{M.\ rosea} DSM 24000$^{\mathrm{T}}$ (C1--C8 attempted; C1--C4 exact; C5 near-exact; C6 exact count with small drift on totals; C7 class-level match; C8 3-of-4 genes confirmed exactly, \emph{pfaE} unconfirmed). C11 arithmetic exact.

\textbf{Not covered:} C9 (ELK/PP ratio across 20 myxobacteria), C10 (HGT-from-Actinobacteria phylogeny), and full domain-scan of PfaA/PfaC/PfaE (PFAM HMMER).

\textbf{LLM-judge consensus:} REPLICATED, score 96/100 (Llama-3.3-70B: 98; Nemotron-3-Ultra: 95; free CELS endpoints only). Note the epistemic caveat in the Critique above.

\textbf{Kept as REPLICATED for the backfill:} yes.

\section{Open Questions}\label{sec:oq}

The five open questions below are grounded in what this replication actually surfaced (see Critique) and what the original paper leaves genuinely unresolved. They are meant to seed follow-on research work.

\begin{description}
\item[Q1.] \textbf{Is \emph{pfaE} really at \texttt{A7982\_13498}, or is the paper's identification annotation-artifactual?} The NCBI CP016211.1 record contains no CDS with a PPTase/Sfp/phosphopantetheinyl-transferase keyword. The paper asserts \emph{pfaE} at \texttt{A7982\_13498} but the underlying GenBank flat file does not label it. Is Fig.\ 5AI's \emph{pfaE} identification an unpublished HMMER call the authors ran but never propagated back into the GenBank submission, or is it a mis-mapping? \emph{Next steps:} run \texttt{hmmscan} of the full \emph{M.\ rosea} proteome against PFAM PF01648 (Sfp/PPTase); check whether \texttt{A7982\_13498}'s translation actually scores as a PPTase; publish a corrected annotation delta if needed.

\item[Q2.] \textbf{Does the AT domain integrated into \emph{pfa3} in \emph{M.\ rosea} vs.\ its absence in \emph{Sorangium} causally explain the DHA/EPA production difference, or is it a correlate?} The paper's central mechanistic hypothesis (Discussion of Fig.\ 5AI vs.\ AIII) is that integrated-AT is why \emph{M.\ rosea} and \emph{Aetherobacter} make DHA/EPA and \emph{Sorangium} does not. This is asserted from comparative genomics, not from biochemistry. \emph{Next steps:} express \emph{M.\ rosea} \emph{pfa1--3} + \emph{pfaE} heterologously in \emph{Streptomyces coelicolor} $\Delta$fatty-acid-null with and without the C-terminal AT domain in \emph{pfa3} deleted; LC-MS profile the resulting fatty acids. This is a lab-scale experiment ($\sim$3--6\,months) but resolves the causal claim.

\item[Q3.] \textbf{Is the ``largest bacterial genome'' claim still true in 2026?} The paper's C11 is time-stamped at 2021 vs.\ \emph{S.\ cellulosum} So0157--2 (14.78\,Mbp). Since 2021 there has been continued PacBio HiFi sequencing of myxobacteria, cyanobacterial \emph{Crocosphaera}, and other large-genome candidates. Neither the paper nor this replication systematically re-surveys the field. \emph{Next steps:} query NCBI \texttt{genome} db for all bacterial assemblies $>$14\,Mbp deposited 2021--2026; produce a ranked table; publish a small update note if the record has been surpassed.

\item[Q4.] \textbf{How much of the ``genome-expansion via gene duplication'' story is actually paralog inflation from repetitive-region assembly artifacts?} The paper reports 44.10\% paralogous genes in \emph{M.\ rosea} vs.\ 47.10\% (\emph{S.\ cellulosum} So ce56) and 41.80\% (\emph{S.\ cellulosum} So0157--2), and uses this to argue duplication drives genome expansion. But PacBio HGAP v2.3.0 (2015-era assembler, used here) can create false duplications at repeat boundaries. The paper does not report a repeat-content analysis (RepeatMasker/RepeatModeler) nor an assembly-QC pass (Merqury, BUSCO duplication rate). \emph{Next steps:} pull the SRA raw reads for CP016211.1, re-assemble with a modern HiFi assembler (hifiasm-meta or Flye 2.9+), compute Merqury duplication metrics, and re-classify the paralog set on the corrected assembly. If duplication rate drops significantly, the paper's central genome-expansion narrative needs rebalancing between ``duplication'' and ``HGT + retention''.

\item[Q5.] \textbf{Is the phylogenetic HGT-from-Actinobacteria call for the \emph{pfa} cluster robust to broader taxonomic sampling and long-branch attraction?} The paper's Fig.\ 5B trees include \emph{Streptomyces}, \emph{Azospirillum}, \emph{Tahibacter}, \emph{Aetherobacter}, \emph{Sorangium} but not the broader $\gamma$-proteobacterial PUFA-PKS carriers (\emph{Shewanella}, \emph{Photobacterium profundum}, \emph{Moritella marina}) as an outgroup, nor a diverse Actinobacterial sweep. The three individual gene trees (Pfa1/Pfa2/Pfa3) may reflect long-branch attraction to \emph{Streptomyces} rather than genuine acquisition. \emph{Next steps:} build a taxon-expanded ($n \sim 100$) concatenated Pfa1+Pfa2+Pfa3 tree using IQ-TREE with ModelFinder + ultrafast bootstrap 1000 + SH-aLRT; test tree topology with an approximately-unbiased (AU) test against the alternate hypothesis of vertical inheritance from a Sorangiineae ancestor; and run an explicit HGT test (RANGER-DTL) with a curated species tree. If AU rejects the Sorangiineae-ancestor topology at $p<0.01$, the HGT claim is strengthened; if not, downgrade C10 to speculation.
\end{description}

\end{document}
